\section{The Relative Zeta Value and an Explicit Upper Bound}\label{an:section}

We bound the relative zeta value $L_F(1)$ needed in the geometric
construction. Logarithms of Dedekind zeta functions are normalized by the
degree: for a number field $A$ we work with $[A:\Q]^{-1}\log\zeta_A(s)$.

\begin{definition}\label{an:family}
Let $\mathcal K$ be the set of all finite Galois extensions $K/B$ contained
in $B_G$, of degree a power of two at least $2^{49}$, such that
$G_B\to\overline G_B$ factors through $\Gal(K/B)$ (notation of
Definition~\ref{tw:GB}). These are the fields $K$ of
Theorem~\ref{tw:field-family}.
\end{definition}

For $K\in\mathcal K$ we use the notation of
Theorem~\ref{tw:field-family}: $\iota_1$ is a complex conjugation at a place of $K$ above
$v_1$, $F=K^{\langle\iota_1\rangle}$, $d=[F:\Q]$, so that $[K:\Q]=2d$, $(b,c)$
is the signature of $F$, $\theta=c/d$, and $\lambda=2^{9/4}\sqrt{3615}$ is the
root discriminant of both $K$ and $F$. Put
\[
 \ell=\log\lambda=\tfrac94\log2+\tfrac12\log3615=5.65600472355\ldots
\]
Different choices of $\iota_1$ are conjugate in $\Gal(K/B)$ and give
conjugate fields $F$; hence $b$, $c$, $\theta$ and the function
$L_F(s)=\zeta_K(s)/\zeta_F(s)$ depend only on $K$. The extension $K/F$ is
quadratic and unramified at every finite place, $K$ is totally imaginary,
and $b\ge1$. Let $\gamma$ be Euler's constant. For a
prime $\mathfrak p$ of $B$, let $e_K(\mathfrak p)$ and $f_K(\mathfrak p)$ be
its ramification index and residue degree in the Galois extension $K/B$, so
that $(e_K(\mathfrak p),f_K(\mathfrak p))$ is the type of $\mathfrak p$ in
$K/B$ (Definition~\ref{tw:type}).

\begin{theorem}\label{an:ceiling}
There is an $m_0$ such that every $K\in\mathcal K$ with $[K:B]\ge m_0$
satisfies
\[
 \frac1d\log L_F(1)<C:=\ceiling.
\]
\end{theorem}

The proof assumes neither the generalized Riemann hypothesis nor a
Brauer--Siegel type asymptotic formula, and it does not give an effective
value of $m_0$. It has three steps. Monotonicity of the completion of $L_F$
reduces the value at $s=1$ to $\zeta_K$ at a real point $\sigma>1$
(Subsection~\ref{an:reduction}). The Tsfasman--Vl\u{a}du\c{t} inequality
bounds the error of this shift (Subsection~\ref{an:budget-section},
Proposition~\ref{an:prime-budget}). Finally, $\zeta_K(\sigma)$ is bounded
prime by prime by $\zeta_E(\sigma)$, for the Kummer field $E$ of
Section~\ref{gf:section}, minus corrections at the primes above $2$, $7$
and $29$, whose types in $K/B$ are known, and at the primes of
$\mathcal P_4$ (Definition~\ref{an:P4}), whose residue degree in $K/B$ is
at least $4$ (Proposition~\ref{an:Ystar}); $\zeta_E(\sigma)$ is evaluated
rigorously in
Subsections~\ref{an:afe-section}--\ref{an:value-section}
(Proposition~\ref{an:genus-value}). Subsection~\ref{an:census} computes
these corrections and the lower bounds for the types in $K/B$ that
Proposition~\ref{an:prime-budget} uses, and the last subsection combines
the steps. Throughout, $\sigma$ denotes a real number greater than $1$ and
$\varepsilon=\sigma-1$; in this section $\varepsilon$ is this positive
number, not the unit $\alpha_1$ of \eqref{tw:alpha}.

\subsection{\texorpdfstring{Reduction to $s>1$}{Reduction to s > 1}}
\label{an:reduction}
For a number field $A$, write $h_A$, $\operatorname{Reg}_A$, $w_A$ and
$\Delta_A$ for its class number, regulator, number of roots of unity and
discriminant. The residue of $\zeta_A$ at $s=1$ is
$2^{r_1(A)}(2\pi)^{r_2(A)}h_A\operatorname{Reg}_A/(w_A|\Delta_A|^{1/2})$.
Since $K/F$ is unramified at the finite places, $|\Delta_K|=|\Delta_F|^2$.
Hence
\begin{equation}\label{an:class-formula}
 L_F(1)=\lim_{s\to1}\frac{\zeta_K(s)}{\zeta_F(s)}
 =2^c\pi^{d-c}\lambda^{-d/2}\,
 \frac{h_K\operatorname{Reg}_Kw_F}{h_F\operatorname{Reg}_Fw_K}>0.
\end{equation}

Let $s>1$ be real. At a prime of $F$ of norm $q$, the Euler factor of $L_F$
is $(1-q^{-s})^{-1}$ if the prime splits in $K$ and $(1+q^{-s})^{-1}$ if it
is inert. Its square is at most the Euler factor of $\zeta_K$, which is
$(1-q^{-s})^{-2}$, respectively $(1-q^{-2s})^{-1}$. Consequently
\begin{equation}\label{an:relative-euler}
 \frac{\log L_F(s)}d\le\frac{\log\zeta_K(s)}{2d}
 =\frac{\log\zeta_K(s)}{[K:\Q]}.
\end{equation}

The function $L_F$ is the Hecke $L$-function of the quadratic character of
$K/F$. This character is unramified at every finite place and nontrivial at
every real place of $F$, since $K$ is totally imaginary. With $\Gamma_\R$,
$\Gamma_\C$ and $\Lambda_A$ as in Section~\ref{gf:section}, put
\begin{equation}\label{an:completion}
 \mathcal L_F(s)=\frac{\Lambda_K(s)}{\Lambda_F(s)}
 =\lambda^{ds/2}\,\Gamma_\R(s+1)^b\,\Gamma_\C(s)^c\,L_F(s).
\end{equation}
By Hecke's theorem \cite{Hecke1920,Tate1967}, $\mathcal L_F$ is entire of
order at most one, and it satisfies $\mathcal L_F(s)=\mathcal L_F(1-s)$
because both completed Dedekind zeta functions do.

\begin{lemma}\label{an:monotonicity}
The function $\mathcal L_F(s)$ is positive and nondecreasing on $[1,\infty)$.
\end{lemma}

\begin{proof}
The Euler product shows that $\mathcal L_F$ has no zeros with
$\operatorname{Re}s>1$, and the functional equation excludes
$\operatorname{Re}s<0$. So every zero $\rho$ satisfies $0\le\operatorname{Re}
\rho\le1$. There is no zero on $[1,\infty)$, by the Euler product and by
\eqref{an:class-formula}. As $\mathcal L_F(s)>0$ for large real $s$,
it is positive on $[1,\infty)$.

Put $\Xi(z)=\mathcal L_F(\frac12+z)$. It is entire of order at most one,
even, and real on the real axis. Its zeros $z=\rho-\frac12$ satisfy
$|\operatorname{Re}z|\le\frac12$, and $\sum_{z\ne0}|z|^{-2}<\infty$. Group
the nonzero zeros into pairs $\{z_0,-z_0\}$, and let $m\ge0$ be the order of
$\Xi$ at $0$. By Hadamard's theorem for functions of order at most one,
\[
 \frac{\Xi'}{\Xi}(z)-\frac mz-\sum_{\{z_0,-z_0\}}\frac{2z}{z^2-z_0^2}
\]
is constant, where the series converges absolutely and locally uniformly
away from the zeros. Since $\Xi$ is even, $\Xi'/\Xi$ is odd, so the constant
is zero. For real $x\ge\frac12$ the logarithmic derivative is real, so it
equals the sum of the real parts. If $z_0=a+iy$, then
\[
 \operatorname{Re}\frac{2x}{x^2-z_0^2}
 =\frac{2x(x^2-a^2+y^2)}{(x^2-a^2+y^2)^2+4a^2y^2}\ge0,
\]
because $|a|\le\frac12\le x$. Since also $m/x\ge0$, the logarithmic
derivative is nonnegative, and $\log\Xi$ is nondecreasing on
$[\frac12,\infty)$.
\end{proof}

For $0<\varepsilon\le1$ and $0\le\theta\le\frac12$ put
\begin{equation}\label{an:gamma}
 \Upsilon(\varepsilon,\theta)
 =\frac\varepsilon2(\ell-\log\pi)-\varepsilon\theta\log2
 +(1-2\theta)\log\Gamma\Bigl(1+\frac\varepsilon2\Bigr)
 +\theta\log\Gamma(1+\varepsilon).
\end{equation}

\begin{lemma}\label{an:shift}
Let $K\in\mathcal K$ and $0<\varepsilon\le1$. Then
\begin{equation}\label{an:finite-gamma}
 \frac{\log L_F(1)}d
 \le\frac{\log\zeta_K(1+\varepsilon)}{[K:\Q]}+\Upsilon(\varepsilon,\theta).
\end{equation}
Moreover $\Upsilon(\varepsilon,\theta)\to0$ as $\varepsilon\downarrow0$,
uniformly for $0\le\theta\le\frac12$.
\end{lemma}

\begin{proof}
By Lemma~\ref{an:monotonicity},
$\mathcal L_F(1)\le\mathcal L_F(1+\varepsilon)$. In \eqref{an:completion}
the gamma factors change by
\[
 \frac{\Gamma_\R(2+\varepsilon)}{\Gamma_\R(2)}
 =\pi^{-\varepsilon/2}\,\Gamma\Bigl(1+\frac\varepsilon2\Bigr),\qquad
 \frac{\Gamma_\C(1+\varepsilon)}{\Gamma_\C(1)}
 =(2\pi)^{-\varepsilon}\,\Gamma(1+\varepsilon).
\]
Taking logarithms, dividing by $d$, and using $b/d=1-2\theta$, $c/d=\theta$
and \eqref{an:relative-euler}, we obtain \eqref{an:finite-gamma}. Since
$\Upsilon$ is affine in $\theta$, $\Upsilon(\varepsilon,\theta)\to0$ as
$\varepsilon\downarrow0$, uniformly for $0\le\theta\le\frac12$.
\end{proof}

The term $\Upsilon$ does not appear in the final estimate. With the value
$\varepsilon=1/300$ used in the proof of Theorem~\ref{an:ceiling},
\eqref{an:finite-gamma} would add $\Upsilon(1/300,\theta)\approx0.0054$ for
$\theta$ near $\frac12$. Instead, the proof of
Proposition~\ref{an:prime-budget} applies \eqref{an:finite-gamma} with
$s-1$ in place of $\varepsilon$ for each fixed $s\in(1,2]$, passes to the
limit along a sequence of fields, and then lets $s\downarrow1$, so that
$\Upsilon$ vanishes in the limit; the
Tsfasman--Vl\u{a}du\c{t} inequality then bounds the change from
$1+\varepsilon$ to $1$ by the term $\varepsilon(\kappa_\infty+\mathcal D)$
of that proposition, which is about $0.00215$ in the proof of
Theorem~\ref{an:ceiling}. This limit is also the reason why $m_0$ is not
effective.

\subsection{\texorpdfstring{The Tsfasman--Vl\u{a}du\c{t} Inequality over $B$}{The Tsfasman-Vladut Inequality over B}}
\label{an:budget-section}
For $q>1$ and $s\ge1$ put $g_s(q)=-\log(1-q^{-s})$.

\begin{lemma}\label{an:local-contribution}
Let $M$ be a finite Galois extension of $B$ and $s>1$. For a prime
$\mathfrak p$ of $B$ let $e_M(\mathfrak p)$ and $f_M(\mathfrak p)$ be its
ramification index and residue degree in $M/B$. Then
\begin{align}
 \frac{\log\zeta_M(s)}{[M:\Q]}
 &=\sum_{\mathfrak p}\frac{g_s(\mathrm N\mathfrak p^{f_M(\mathfrak p)})}
 {2e_M(\mathfrak p)f_M(\mathfrak p)},\label{an:local-sum}\\
 -\frac1{[M:\Q]}\frac{\zeta_M'}{\zeta_M}(2)
 &=\frac1{[M:\Q]}\sum_{\mathfrak P}\frac{\log\mathrm N\mathfrak P}
 {\mathrm N\mathfrak P^2-1}
 =\sum_{\mathfrak p}\frac{\log\mathrm N\mathfrak p}
 {2e_M(\mathfrak p)(\mathrm N\mathfrak p^{2f_M(\mathfrak p)}-1)},
 \label{an:local-debit}
\end{align}
where $\mathfrak P$ runs over the primes of $M$. Each summand on the right
does not increase when $e_M(\mathfrak p)$ or $f_M(\mathfrak p)$ is replaced
by a larger real number. If $M\subseteq M'$ are both Galois over $B$, then
$e_M(\mathfrak p)\mid e_{M'}(\mathfrak p)$ and
$f_M(\mathfrak p)\mid f_{M'}(\mathfrak p)$.
\end{lemma}

\begin{proof}
Above $\mathfrak p$ there are $[M:B]/(ef)$ primes of $M$, each of norm
$\mathrm N\mathfrak p^f$, and $[M:\Q]=2[M:B]$. This gives both formulas.
The first equality in \eqref{an:local-debit} follows from
$-\zeta_M'/\zeta_M(s)=\sum_{\mathfrak P}\log\mathrm N\mathfrak P/(\mathrm
N\mathfrak P^s-1)$. For monotonicity in $f$, write
$g_s(x^f)/f=\sum_{j\ge1}x^{-fjs}/(fj)$ for $x>1$; each term decreases in
$f$. The divisibilities hold because ramification indices and residue
degrees are multiplicative in towers.
\end{proof}

Put
\[
 \kappa_\infty=\frac{\ell-\gamma-\log(4\pi)}4,\qquad
 w(q)=2\log q\sum_{m\ge1}\frac1{q^m+1}.
\]
Here $2\kappa_\infty$ is the right side, and $w(q)$ the weight of the prime
power $q$, in the Tsfasman--Vl\u{a}du\c{t} inequality \eqref{an:tv-budget}
below.

\begin{proposition}\label{an:prime-budget}
Let $\sigma=1+\varepsilon$ with $0<\varepsilon\le1$, and let $Y_*$ be a real
number with $[K:\Q]^{-1}\log\zeta_K(\sigma)\le Y_*$ for every
$K\in\mathcal K$. For each prime $\mathfrak p$ of $B$, let
$e^0_{\mathfrak p},f^0_{\mathfrak p}\ge1$ be integers with
$e_K(\mathfrak p)\ge e^0_{\mathfrak p}$ and $f_K(\mathfrak p)\ge
f^0_{\mathfrak p}$ for every $K\in\mathcal K$, that is, lower bounds for
the type of $\mathfrak p$ in $K/B$, and let
\[
 \mathcal D\ge\sum_{\mathfrak p}\frac{\log\mathrm N\mathfrak p}
 {2e^0_{\mathfrak p}(\mathrm N\mathfrak p^{2f^0_{\mathfrak p}}-1)}.
\]
Then every $C'>Y_*+\varepsilon(\kappa_\infty+\mathcal D)$ has the following
property: there is an $m_0$ such that $d^{-1}\log L_F(1)<C'$ for every
$K\in\mathcal K$ with $[K:B]\ge m_0$.
\end{proposition}

\begin{proof}
\emph{A limiting sequence.} Suppose not. Then there are fields $K^{(j)}\in\mathcal K$ of strictly
increasing degree, hence pairwise non-isomorphic, with
$d_j^{-1}\log L_{F^{(j)}}(1)\ge C'$, where $[K^{(j)}:\Q]=2d_j$. For a prime power $q$
let $N_q(K)$ be the number of primes of $K$ of norm $q$. If $q=p^f$, then
$N_q(K)\le[K:\Q]/f$, since the local degrees above $p$ add up to $[K:\Q]$.
By a diagonal argument we may assume that
$\beta_q=\lim_jN_q(K^{(j)})/(2d_j)$ exists for every $q$.

\emph{The Tsfasman--Vl\u{a}du\c{t} inequality.} Tsfasman and
Vl\u{a}du\c{t} call a sequence $(A_i)$ of pairwise
non-isomorphic number fields asymptotically exact if its genus
$\mathfrak g(A_i)=\log|\Delta_{A_i}|^{1/2}$ tends to infinity and the limits
$\phi_q=\lim N_q(A_i)/\mathfrak g(A_i)$,
$\phi_\R=\lim r_1(A_i)/\mathfrak g(A_i)$ and
$\phi_\C=\lim r_2(A_i)/\mathfrak g(A_i)$ exist. Their unconditional Basic
Inequality$'$ \cite[Section~3.2, Proposition~3.1]{TsfasmanVladut2002} states
that every such sequence satisfies
\[
 2\sum_q\phi_q\log q\sum_{m\ge1}\frac1{q^m+1}
 +\phi_\R\Bigl(\frac\gamma2+\frac12+\log2\sqrt\pi\Bigr)
 +\phi_\C(\gamma+\log4\pi)\le1.
\]
For our sequence $\mathfrak g(K^{(j)})=d_j\ell$ and $K^{(j)}$ is totally imaginary, so
$\phi_q=2\beta_q/\ell$, $\phi_\R=0$ and $\phi_\C=1/\ell$, and the inequality
becomes
\begin{equation}\label{an:tv-budget}
 \sum_q\beta_qw(q)\le\frac{\ell-\gamma-\log(4\pi)}2=2\kappa_\infty.
\end{equation}

\emph{The value at $1$.} Put $Z(s)=\sum_q\beta_qg_s(q)$ for $s\ge1$. Since $g_1(q)\le1/(q-1)$ and
$w(q)\ge2\log q/(q+1)$, we have $g_1(q)\le3w(q)/(2\log2)$, so $Z(1)$ is
finite. For fixed real $s>1$,
$(2d_j)^{-1}\log\zeta_{K^{(j)}}(s)=\sum_q(N_q(K^{(j)})/2d_j)g_s(q)$ tends to $Z(s)$ by
dominated convergence: for $q=p^f$ the terms are at most
$g_s(q)/f\le2q^{-s}$, and $\sum_q q^{-s}<\infty$. By Lemma~\ref{an:shift}
with $s-1$ in place of $\varepsilon$, for $1<s\le2$,
\[
 C'\le\limsup_j\frac{\log L_{F^{(j)}}(1)}{d_j}
 \le Z(s)+\sup_{0\le\theta\le1/2}\Upsilon(s-1,\theta).
\]
As $s\downarrow1$, $Z(s)$ increases to $Z(1)$ and the supremum tends to
zero. Hence $C'\le Z(1)$.

\emph{From $1+\varepsilon$ to $1$.} For every prime power $q$,
\begin{align*}
 g_1(q)-g_{1+\varepsilon}(q)&=\int_1^{1+\varepsilon}\frac{\log q}{q^u-1}\,du
 \le\frac{\varepsilon\log q}{q-1},\\
 \frac{\log q}{q-1}-\frac{w(q)}2&=\log q\sum_{m\ge1}\frac1{q^m(q^m+1)}
 \le\frac{\log q}{q^2-1}.
\end{align*}
With \eqref{an:tv-budget} these give
\[
 Z(1)\le Z(1+\varepsilon)
 +\varepsilon\Bigl(\kappa_\infty+\sum_q\beta_q\frac{\log q}{q^2-1}\Bigr).
\]
By Fatou's lemma, \eqref{an:local-debit} for $M=K^{(j)}$, and the monotonicity
in Lemma~\ref{an:local-contribution},
\[
 \sum_q\beta_q\frac{\log q}{q^2-1}
 \le\liminf_j\sum_{\mathfrak p}\frac{\log\mathrm N\mathfrak p}
 {2e_{K^{(j)}}(\mathfrak p)(\mathrm N\mathfrak p^{2f_{K^{(j)}}(\mathfrak p)}-1)}
 \le\mathcal D.
\]
Finally
$Z(1+\varepsilon)=\lim_j(2d_j)^{-1}\log\zeta_{K^{(j)}}(1+\varepsilon)\le Y_*$.
Thus $C'\le Z(1)\le Y_*+\varepsilon(\kappa_\infty+\mathcal D)$, a
contradiction.
\end{proof}

\subsection{Comparison with the Kummer Field}\label{an:comparison}
We use the set $S$ of the six primes of $B$ above $2$, $3$ and $5$ and the
primes $\mathfrak t_0=7\mathcal O_B$ and $\mathfrak t_1,\mathfrak t_2$ above
$29$ (Lemma~\ref{tw:base}), the group $G_B$ of Definition~\ref{tw:GB} and
its Zassenhaus filtration $D_nG_B$, the relation space $R_2$ in degree two
(Definition~\ref{tw:relation-space}), which Lemma~\ref{tw:quadratic-layer}
computes, the restricted square $v^{[2]}$ of a vector $v\in\Ftwo^8$ as in
Section~\ref{tw:tower}, and the Frobenius vectors
$\overline{\Frob}_{\mathfrak p}$ of Definition~\ref{gf:frobenius-vector}.

\begin{definition}\label{an:P4}
Let $\Sigma_2=\{v\in\Ftwo^8:v^{[2]}\in R_2\}$, and let $\mathcal P_4$ be the
set of primes $\mathfrak p\notin S\cup\{\mathfrak t_0,\mathfrak
t_1,\mathfrak t_2\}$ of $B$ with $\mathrm N\mathfrak p\le10^6$ and
$\overline{\Frob}_{\mathfrak p}\notin\Sigma_2$.
\end{definition}

\begin{lemma}\label{an:census-degree}
Let $K\in\mathcal K$, and let $\mathfrak p\notin S$ be a prime of $B$ with
$\overline{\Frob}_{\mathfrak p}\notin\Sigma_2$. Then $e_K(\mathfrak p)=1$ and
$f_K(\mathfrak p)\ge4$.
\end{lemma}

\begin{proof}
The field $K$ is unramified over $B$ outside $S$, so $e_K(\mathfrak p)=1$,
and a prime $\mathfrak P$ of $K$ above $\mathfrak p$ has a Frobenius
$\Frob_{\mathfrak P}\in\Gal(K/B)$ of order $f_K(\mathfrak p)$. Choose
$g\in G_B$ mapping to $\Frob_{\mathfrak P}$. Its image in
$G_B/D_2G_B=\Gal(E/B)$ is the Frobenius of $\mathfrak p$ in $E$
(Lemma~\ref{gf:in-tower}), whose elementary image is $\overline{\Frob}_{\mathfrak p}$; it is
nonzero, since $0\in\Sigma_2$. By Lemma~\ref{tw:quadratic-layer}, which uses
the complete list of generators of $N_B$ in
Lemma~\ref{tw:GB-presentation}, and since $\overline{\Frob}_{\mathfrak p}^{[2]}\notin R_2$,
$g$ has order at least $4$ in $G_B/D_3G_B$. The projection $G_B\to G_B/D_3G_B$
factors through $\Gal(K/B)$, because $K$ contains the fixed field of
$D_4G_B\subseteq D_3G_B$. So $\Frob_{\mathfrak P}$ has order at least $4$.
\end{proof}

The next proposition compares $\zeta_K(\sigma)$ with $\zeta_E(\sigma)$ and
subtracts two corrections. The first concerns the primes
$\mathfrak p_1,\mathfrak p_2$ above $2$ and $\mathfrak t_0,\mathfrak
t_1,\mathfrak t_2$ above $7$ and $29$, whose types in $K/B$ are known; the
second concerns the primes of $\mathcal P_4$, which are found in
Subsection~\ref{an:census}. Put
\begin{align}
 \Delta_{2,7,29}(\sigma)&=2\Bigl(\frac{g_\sigma(2^2)}{16}
 -\frac{g_\sigma(2^4)}{64}\Bigr)
 +2\Bigl(\frac{g_\sigma(29^2)}4-\frac{g_\sigma(29^4)}8\Bigr)
 +\frac{g_\sigma(7^4)}4-\frac{g_\sigma(7^8)}8,
 \label{an:sel}\\
 \Delta_{\mathcal P_4}(\sigma)&=\sum_{\mathfrak p\in\mathcal P_4}
 \Bigl(\frac{g_\sigma(\mathrm N\mathfrak p^2)}4
 -\frac{g_\sigma(\mathrm N\mathfrak p^4)}8\Bigr).
 \label{an:cen}
\end{align}
Every term of both sums is nonnegative, since $g_\sigma(x^2)\le g_\sigma(x)$.

\begin{proposition}\label{an:Ystar}
Every $K\in\mathcal K$ satisfies
\[
 \frac{\log\zeta_K(\sigma)}{[K:\Q]}
 \le\frac{\log\zeta_E(\sigma)}{512}-\Delta_{2,7,29}(\sigma)
 -\Delta_{\mathcal P_4}(\sigma).
\]
\end{proposition}

\begin{proof}
By Lemma~\ref{gf:in-tower}, $E\subseteq K$, and both are Galois over $B$.
Apply \eqref{an:local-sum} to $M=E$ and $M=K$. By
Lemma~\ref{an:local-contribution}, each term for $K$ is at most the
corresponding term for $E$. For the primes listed below we bound the term
for $K$ from above using what is known about the type in $K/B$, and
subtract the difference between the term for $E$ and this upper bound.

The types in $E/B$ are given by Corollary~\ref{gf:types}: $(4,2)$ at
$\mathfrak p_1,\mathfrak p_2$, and $(1,2)$ at $\mathfrak t_0,\mathfrak
t_1,\mathfrak t_2$ and at the primes of $\mathcal P_4$. By
Theorem~\ref{tw:field-family}, the primes of $K$ above $2$, $29$ and $7$ have
absolute types $(8,4)$, $(1,4)$ and $(1,8)$. The primes $\mathfrak p_j$,
$\mathfrak t_1$ and $\mathfrak t_2$ have residue degree one over $\Q$, and
$\mathfrak t_0$ has residue degree two. So the types in $K/B$ are $(8,4)$ at
$\mathfrak p_j$ and $(1,4)$ at $\mathfrak t_0,\mathfrak t_1,\mathfrak t_2$.
By Lemma~\ref{an:census-degree}, the primes of $\mathcal P_4$ have $e_K=1$
and $f_K\ge4$. The differences of the terms of $E$ and the upper bounds for
the terms of $K$ are exactly the summands of \eqref{an:sel} and
\eqref{an:cen}. For example, at $\mathfrak p_j$ the term of $E$ is
$g_\sigma(2^2)/(2\cdot4\cdot2)$ and that of $K$ is
$g_\sigma(2^4)/(2\cdot8\cdot4)$.
\end{proof}

\subsection{An Approximate Functional Equation for Two Gamma Factors}
\label{an:afe-section}
We evaluate $\zeta_E(\sigma)$ through Proposition~\ref{gf:factorization}. The
factors $L(\sigma,\chi_e)$ are computed from an approximate functional
equation. For real $\mu$ and $x,y>0$ define the incomplete gamma and Bessel
integrals
\[
 H_\mu(x)=\int_1^\infty u^{\mu-1}e^{-xu}\,du,\qquad
 I_\mu(y)=\int_1^\infty u^\mu K_0(yu)\,du,
\]
where $K_0$ is the modified Bessel function of the second kind of order
zero, and their envelopes (upper bounds, by Lemma~\ref{an:kernels}(2))
\[
 \mathcal V_1(x)=\frac{e^{-x}}x\Bigl(1+\frac1x\Bigr),\qquad
 \mathcal V_2(y)=\sqrt{\frac\pi2}\,\frac{e^{-y}}{y^{3/2}}\Bigl(1+\frac1y\Bigr).
\]
Numerical evaluation of $L$-functions from their functional equations is
developed in \cite{Dokchitser2004}. Part (1) of the next lemma treats mixed
gamma factors, and part (2), for pure gamma factors, uses the Bessel
function $K_0$. We prove both.

\begin{lemma}\label{an:balanced-afe}
Let $a_n$ be real numbers with $|a_n|\le\tau(n)$, let $Q>0$ and
$\nu_1,\nu_2\in\{0,1\}$ (integers, not places), and let
$L(s)=\sum_{n\ge1}a_nn^{-s}$. Suppose that
$\Lambda(s)=Q^{s/2}\Gamma_\R(s+\nu_1)\Gamma_\R(s+\nu_2)L(s)$ extends to an
entire function of order at most one with $\Lambda(1-s)=\Lambda(s)$. Put
$t=2\pi/\sqrt Q$ and let $\sigma>1$ be real.
\begin{enumerate}
\item If $\nu_1\ne\nu_2$, then
\[
 L(\sigma)=\frac{t^\sigma}{\Gamma(\sigma)}\sum_{n\ge1}a_n
 \bigl(H_\sigma(tn)+H_{1-\sigma}(tn)\bigr).
\]
\item If $\nu_1=\nu_2=\nu$, then
\[
 L(\sigma)=\frac{4(t/2)^{\sigma+\nu}}{\Gamma((\sigma+\nu)/2)^2}
 \sum_{n\ge1}a_nn^\nu
 \bigl(I_{\sigma+\nu-1}(tn)+I_{\nu-\sigma}(tn)\bigr).
\]
\end{enumerate}
The series converge absolutely.
\end{lemma}

\begin{proof}
In case (1) put $k(x)=e^{-x}$ and $\nu=0$, and in case (2) put
$k(x)=K_0(x)$. Their Mellin transforms are $\Gamma(w)$ and
$2^{w-2}\Gamma(w/2)^2$ for $\operatorname{Re}w>0$
\cite[10.43.19]{DLMF}. Let $\Theta(u)=\sum_na_nn^\nu k(tnu)$ for $u>0$. For
$\operatorname{Re}w>1+\nu$, termwise integration gives
\[
 \widetilde\Theta(w)=\int_0^\infty\Theta(u)u^{w-1}\,du=
 \begin{cases}
 t^{-w}\Gamma(w)L(w)=\tfrac12\Lambda(w)&\text{in case (1)},\\
 2^{w-2}t^{-w}\Gamma(w/2)^2L(w-\nu)=\tfrac14Q^{\nu/2}\Lambda(w-\nu)
 &\text{in case (2)}.
 \end{cases}
\]
In case (1) we used $\Gamma_\R(w)\Gamma_\R(w+1)=\Gamma_\C(w)$ and
$Q^{w/2}(2\pi)^{-w}=t^{-w}$. In both cases $\widetilde\Theta$ is entire and
$\widetilde\Theta(w)=\widetilde\Theta(1+2\nu-w)$.

The function $L=\Lambda/(Q^{s/2}\Gamma_\R(s+\nu_1)\Gamma_\R(s+\nu_2))$ of the
complex variable $s$ is entire of finite order. Fix $\eta>0$. The function
$L$ is bounded on $\operatorname{Re}s=1+\eta$, and, by the functional
equation and Stirling's formula, of polynomial growth on
$\operatorname{Re}s=-\eta$. By the Phragm\'en--Lindel\"of principle it has
polynomial growth in $|\operatorname{Im}s|$ on every vertical strip. With
Stirling's formula this shows that $\widetilde\Theta(w)$ decays exponentially in
$|\operatorname{Im}w|$, uniformly on vertical strips. Mellin inversion on a
line $\operatorname{Re}w=c>1+\nu$ and a shift to
$\operatorname{Re}w=1+2\nu-c$ are therefore justified, and no residues
occur. With the functional equation of $\widetilde\Theta$ they give
\[
 \Theta(u)=u^{-1-2\nu}\,\Theta(1/u)\qquad(u>0).
\]
In particular $\Theta(u)$ decays rapidly as $u\downarrow0$. Splitting the
Mellin integral at $u=1$ and substituting $u\mapsto1/u$ below $1$ gives
\[
 \widetilde\Theta(\sigma+\nu)=\int_1^\infty\Theta(u)
 \bigl(u^{\sigma+\nu-1}+u^{\nu-\sigma}\bigr)\,du.
\]
Termwise integration, justified by the exponential decay of $k$, turns the
right side into $\sum_na_nn^\nu$ times $H_\sigma(tn)+H_{1-\sigma}(tn)$ in case
(1), and times $I_{\sigma+\nu-1}(tn)+I_{\nu-\sigma}(tn)$ in case (2). On the
left, $\widetilde\Theta(\sigma)=t^{-\sigma}\Gamma(\sigma)L(\sigma)$ in case (1), and
$\widetilde\Theta(\sigma+\nu)=2^{\sigma+\nu-2}t^{-\sigma-\nu}
\Gamma((\sigma+\nu)/2)^2L(\sigma)$ in case (2).
\end{proof}

\begin{lemma}\label{an:kernels}
Let $\mu\in\R$.
\begin{enumerate}
\item The functions $H_\mu$ and $I_\mu$ are Laplace transforms of positive
 measures on $[1,\infty)$. They are completely monotone on $(0,\infty)$,
 extend holomorphically to $\operatorname{Re}z>0$, and satisfy
 $|H(z)|\le H(\operatorname{Re}z)$ there, for $H=H_\mu$ or $H=I_\mu$.
\item If $\mu\le2$, then $H_\mu\le\mathcal V_1$. If $\mu\le3/2$, then
 $I_\mu\le\mathcal V_2$.
\item They satisfy
 \[
 xH_\mu'(x)+\mu H_\mu(x)=-e^{-x},\qquad
 yI_\mu'(y)+(\mu+1)I_\mu(y)=-K_0(y),
 \]
 and $K_0$ satisfies $y^2K_0''+yK_0'-y^2K_0=0$ and $K_0'=-K_1$, where
 $K_1$ is the modified Bessel function of the second kind of order one.
\end{enumerate}
\end{lemma}

\begin{proof}
The definition exhibits $H_\mu$ as the Laplace transform of
$u^{\mu-1}\,du$ on $[1,\infty)$. Taking the order $0$ in
\cite[10.32.9]{DLMF} and substituting $v=\cosh r$ in its integral gives
$K_0(x)=\int_1^\infty e^{-xv}(v^2-1)^{-1/2}\,dv$. Substituting this into
$I_\mu$, and then $v=r/u$ in the inner integral, shows
\[
 I_\mu(y)=\int_1^\infty e^{-yr}\Bigl(\int_1^ru^\mu(r^2-u^2)^{-1/2}\,du\Bigr)
 \,dr,
\]
with a nonnegative inner integral. Differentiation under the integral
gives complete monotonicity. Since the measures are positive and
$|e^{-zr}|=e^{-r\operatorname{Re}z}$, we get $|H(z)|\le H(\operatorname{Re}z)$
for $\operatorname{Re}z>0$, which proves (1).

For (2), $u^{\mu-1}\le u$ for $u\ge1$ and $\mu\le2$, and
$\int_1^\infty ue^{-xu}\,du=\mathcal V_1(x)$. For $I_\mu$ we use
$K_0(x)\le K_{1/2}(x)=\sqrt{\pi/(2x)}\,e^{-x}$ \cite[10.39.2]{DLMF}. It
holds because the modified Bessel function of order $\eta\ge0$ is
$\int_0^\infty e^{-x\cosh r}\cosh(\eta r)\,dr$ \cite[10.32.9]{DLMF}, which
increases with $\eta$. Since $u^{\mu-1/2}\le u$ for $\mu\le3/2$,
\[
 I_\mu(y)\le\sqrt{\frac\pi{2y}}\int_1^\infty u^{\mu-1/2}e^{-yu}\,du
 \le\mathcal V_2(y).
\]

For (3), differentiate under the integral and integrate by parts:
$xH_\mu'(x)=-\int_1^\infty u^\mu xe^{-xu}\,du=-e^{-x}-\mu H_\mu(x)$, and
$yI_\mu'(y)=\int_1^\infty u^{\mu+1}\frac{d}{du}K_0(yu)\,du
=-K_0(y)-(\mu+1)I_\mu(y)$. The last two identities are Bessel's equation
and a standard derivative \cite[10.25.1, 10.29.3]{DLMF}.
\end{proof}

In the application $\sigma=301/300$, so $1<\sigma\le3/2$. The four exponents $\sigma$,
$1-\sigma$, $\sigma+\nu-1$ and $\nu-\sigma$ in Lemma~\ref{an:balanced-afe}
are then covered by the envelopes: $\mathcal V_1$ bounds $H_\sigma$ and
$H_{1-\sigma}$, and $\mathcal V_2$ bounds $I_{\sigma+\nu-1}$ and
$I_{\nu-\sigma}$.

\begin{lemma}\label{an:tails}
Let $1<\sigma\le3/2$, let $|a_n|\le\tau(n)$, let $t>0$, $N\ge1$ and
$\omega>1$.
\begin{enumerate}
\item If $t(N+1)\ge\omega-1$, then for $\mu\in\{\sigma,1-\sigma\}$
\[
 \sum_{n>N}|a_n|H_\mu(tn)\le\zeta(\omega)^2(N+1)^\omega
 \mathcal V_1\bigl(t(N+1)\bigr).
\]
\item If $\nu\in\{0,1\}$ and $t(N+1)\ge\omega+\nu-\frac32$, then for
 $\mu\in\{\sigma+\nu-1,\nu-\sigma\}$
\[
 \sum_{n>N}|a_n|n^\nu I_\mu(tn)\le\zeta(\omega)^2(N+1)^{\omega+\nu}
 \mathcal V_2\bigl(t(N+1)\bigr).
\]
\end{enumerate}
\end{lemma}

\begin{proof}
The logarithmic derivative of
$x^\omega\mathcal V_1(x)=x^{\omega-1}e^{-x}(1+1/x)$ is
$(\omega-1)/x-1-1/(x(x+1))<0$ for $x\ge\omega-1$. So
$n^\omega\mathcal V_1(tn)\le(N+1)^\omega\mathcal V_1(t(N+1))$ for $n>N$. By
Lemma~\ref{an:kernels}(2),
\[
 \sum_{n>N}|a_n|H_\mu(tn)\le\sum_{n>N}\frac{\tau(n)}{n^\omega}\,
 n^\omega\mathcal V_1(tn)\le(N+1)^\omega\mathcal V_1\bigl(t(N+1)\bigr)
 \sum_{n\ge1}\frac{\tau(n)}{n^\omega},
\]
and the last sum is $\zeta(\omega)^2$. Part (2) is the same, since
$y^{\omega+\nu}\mathcal V_2(y)$ is a multiple of
$y^{\omega+\nu-3/2}e^{-y}(1+1/y)$, which decreases for
$y\ge\omega+\nu-\frac32$.
\end{proof}

\subsection{\texorpdfstring{Enclosures of $H_\mu$ and $I_\mu$, and Finite Sums}{Enclosures of H and I, and Finite Sums}}
\label{an:kernel-enclosures}
An \emph{enclosure} of a real number is an interval that contains it. We
enclose $H_\mu$ and $I_\mu$ on a geometric grid. Let $H$ be $H_\mu$ or $I_\mu$
with $\mu$ as in Lemma~\ref{an:tails}, and let $\mathcal V$ be its envelope,
$\mathcal V_1$ or $\mathcal V_2$. For $J\ge0$ and $\xi>0$ put
\[
 \rho_J(\xi)=\mathcal V(\xi/2)\,\frac{16^{-J-1}}{1-1/16}.
\]

\begin{lemma}\label{an:taylor}
Let $\xi>0$ and let $h_j=H^{(j)}(\xi)/j!$ be the Taylor coefficients of $H$
at $\xi$.
\begin{enumerate}
\item $(-1)^jh_j\ge0$ and $|h_j|\le\mathcal V(\xi/2)(2/\xi)^j$ for all
 $j\ge0$.
\item For $|x-\xi|\le\xi/32$,
 $\bigl|H(x)-\sum_{j=0}^Jh_j(x-\xi)^j\bigr|\le\rho_J(\xi)$. For
 $x=31\xi/32$ the difference lies in $[0,\rho_J(\xi)]$.
\item Let $k_j$ be the Taylor coefficients at $\xi$ of $e^{-x}$ if
 $H=H_\mu$, and of $K_0$ if $H=I_\mu$. If $H=H_\mu$, then
 $(j+1)\xi h_{j+1}=-k_j-(j+\mu)h_j$, and $k_j=(-1)^je^{-\xi}/j!$. If
 $H=I_\mu$, then $(j+1)\xi h_{j+1}=-k_j-(j+\mu+1)h_j$, where
 $k_0=K_0(\xi)$, $k_1=-K_1(\xi)$, and
 \[
 (j+1)(j+2)\xi^2k_{j+2}=-(j+1)(2j+1)\xi k_{j+1}-(j^2-\xi^2)k_j
 +2\xi k_{j-1}+k_{j-2},
 \]
 with $k_{-1}=k_{-2}=0$.
\end{enumerate}
\end{lemma}

\begin{proof}
Complete monotonicity gives the signs in (1). The disc $|z-\xi|\le\xi/2$ lies
in $\operatorname{Re}z\ge\xi/2$, where $|H(z)|\le H(\xi/2)\le\mathcal
V(\xi/2)$ by Lemma~\ref{an:kernels}. Cauchy's estimate gives the bound on
$h_j$. For $|x-\xi|\le\xi/32$, the omitted terms have modulus at most
$\mathcal V(\xi/2)16^{-j}$, and their sum over $j>J$ is at most
$\rho_J(\xi)$. For $x<\xi$ every term $h_j(x-\xi)^j$ is nonnegative by (1),
so the remainder is nonnegative. Part (3) follows by comparing coefficients
of $(x-\xi)^j$ in the differential equations of Lemma~\ref{an:kernels}(3),
written at $x=\xi+(x-\xi)$.
\end{proof}

The grid points are $\xi_i=64\,(31/32)^i$ for $0\le i\le i_1$, where
$i_1=1878$ is the least index with $\xi_{i_1}\le2^{-80}$. Taylor polynomials
are formed at $\xi_0,\ldots,\xi_{i_1-1}$. The initial values
$e^{-\xi_i}$, $K_0(\xi_i)$ and $K_1(\xi_i)$ of part~(3) of
Lemma~\ref{an:taylor} are enclosed with the rigorous exponential and Bessel
functions of Arb \cite{Johansson2017}, which is now part of FLINT
\cite{FLINT} and is called through its Python interface.

\begin{lemma}\label{an:grid}
Let $J\ge0$, and suppose that enclosures of $e^{-\xi_i}$, $K_0(\xi_i)$ and
$K_1(\xi_i)$ are given for $0\le i<i_1$. Start from the enclosure
$[0,\mathcal V(64)]$ of $H(\xi_0)$. For $i=0,1,\ldots,i_1-1$, compute from
the enclosure of $H(\xi_i)$ enclosures of $h_0,\ldots,h_J$ at $\xi_i$ by the
recurrences of Lemma~\ref{an:taylor}(3), and take
\[
 \sum_{j=0}^Jh_j(-\xi_i/32)^j+[0,\rho_J(\xi_i)],
\]
evaluated in interval arithmetic, as the enclosure of $H(\xi_{i+1})$. Then
every interval so obtained contains the corresponding true value.
\end{lemma}

\begin{proof}
By Lemma~\ref{an:kernels}, $H(\xi_0)\in[0,\mathcal V(64)]$. Suppose that
the enclosure of $H(\xi_i)$ contains $H(\xi_i)$. The true Taylor
coefficients at $\xi_i$ are obtained from the true value $H(\xi_i)$ by the
same recurrence, so their enclosures contain them. Since
$\xi_{i+1}=31\xi_i/32$, Lemma~\ref{an:taylor}(2) shows that the enclosure
of $H(\xi_{i+1})$ contains $H(\xi_{i+1})$. Induction on $i$ proves the
lemma.
\end{proof}

Now let $t>0$, $\nu\in\{0,1\}$, integers $a_1,\ldots,a_N$ and $J\ge0$ be
given, with $tN<\xi_0$ and $t>\xi_{i_1}$. For $0\le i<i_1$ let
$\mathcal B_i=\{n:1\le n\le N,\ \xi_{i+1}<tn\le\xi_i\}$. These sets
partition $\{1,\ldots,N\}$. For nonempty $\mathcal B_i$, let $\bar n_i$ be
the integer part of the average of its least and largest elements, and put
\[
 M_{i,j}=\sum_{n\in\mathcal B_i}a_nn^\nu(n-\bar n_i)^j,\qquad
 A_i=\sum_{n\in\mathcal B_i}|a_n|n^\nu,\qquad
 P_{i,j}=t^j\sum_{m=j}^J\binom mjh_{i,m}(t\bar n_i-\xi_i)^{m-j},
\]
where $h_{i,m}$ are the Taylor coefficients of $H$ at $\xi_i$. The moments
$M_{i,j}$ and the absolute sums $A_i$ are exact integers.

\begin{lemma}\label{an:binning}
With this notation,
\begin{equation}\label{an:bin-sum}
 \Bigl|\sum_{n=1}^Na_nn^\nu H(tn)-\sum_i\sum_{j=0}^JP_{i,j}M_{i,j}\Bigr|
 \le\sum_i\rho_J(\xi_i)A_i,
\end{equation}
where $i$ runs over the indices with $\mathcal B_i$ nonempty.
\end{lemma}

\begin{proof}
Let $n\in\mathcal B_i$. Then $|tn-\xi_i|<\xi_i/32$, so by
Lemma~\ref{an:taylor}(2) the value $H(tn)$ differs from
$\sum_{m=0}^Jh_{i,m}(tn-\xi_i)^m$ by at most $\rho_J(\xi_i)$. Since
$tn-\xi_i=t(n-\bar n_i)+(t\bar n_i-\xi_i)$, the binomial theorem gives
$\sum_{m=0}^Jh_{i,m}(tn-\xi_i)^m=\sum_{j=0}^JP_{i,j}(n-\bar n_i)^j$.
Multiply by $a_nn^\nu$ and sum over $n\in\mathcal B_i$ and over $i$.
\end{proof}

In the application below $t\ge2\pi/\sqrt{3470400}>0.0033$, because
$Q_e\le3470400$ for every $e$ (Corollary~\ref{gf:conductors}). So every argument
$tn$ exceeds $0.0033$, and the part of the grid below that point is not
used.

\subsection{\texorpdfstring{The Value of $\zeta_E$ at $\sigma=301/300$}{The Value of zeta(E) at sigma = 301/300}}
\label{an:value-section}
From now on $\sigma=301/300$. Let $e\in\Ftwo^8\setminus\{0\}$. By
Proposition~\ref{gf:degree-two}, $L(s,\chi_e)$ satisfies the hypotheses of
Lemma~\ref{an:balanced-afe} with $Q=Q_e$ and $\nu_k=\nu_k(e)$. Put
$t=2\pi/\sqrt{Q_e}$ and
\[
 N_e=4\lfloor\sqrt{Q_e}\rfloor+1,
\]
and let $\Pi_e$ be the prefactor of that lemma:
$\Pi_e=t^\sigma/\Gamma(\sigma)$ if $\nu_1(e)\ne\nu_2(e)$, and
$\Pi_e=4(t/2)^{\sigma+\nu}/\Gamma((\sigma+\nu)/2)^2$ if
$\nu_1(e)=\nu_2(e)=\nu$; in the first case put $\nu=0$. Let
$\mathcal S_1$ and $\mathcal S_2$ be the sums
$\sum_i\sum_{j=0}^{20}P_{i,j}M_{i,j}$ of Lemma~\ref{an:binning}, with
$N=N_e$, $a_n=a_n(e)$ and $J=20$, for the two functions of the series of
Lemma~\ref{an:balanced-afe}: $H=H_\sigma$ and $H=H_{1-\sigma}$ in the first
case, and $H=I_{\sigma+\nu-1}$ and $H=I_{\nu-\sigma}$ in the second; their
envelope is $\mathcal V=\mathcal V_1$ in the first case and
$\mathcal V=\mathcal V_2$ in the second. For every $e$, $N_e$ lies between
$105$ and $7449$, and $t(N_e+1)$ between $24.7$ and $25.5$ (the
supplementary program \texttt{check255.gp} prints both ranges). Hence
$tN_e<\xi_0$ and $t>\xi_{i_1}$, as Lemma~\ref{an:binning} requires, and the
hypotheses of Lemma~\ref{an:tails} hold for every $\omega$ in the list
$\frac{101}{100},\frac{21}{20},\frac{11}{10},\frac98,\frac65,\frac54,
\frac43,\frac32,2$. The value of $\omega$ in this list is chosen to minimize
the bound
\[
 \mathcal T_e=\sum_i\rho_{20}(\xi_i)A_i
 +\zeta(\omega)^2(N_e+1)^{\omega+\nu}\mathcal V\bigl(t(N_e+1)\bigr).
\]

\begin{lemma}\label{an:row}
For $\sigma=301/300$ and every $e\in\Ftwo^8\setminus\{0\}$,
\begin{equation}\label{an:row-enclosure}
 \bigl|L(\sigma,\chi_e)-\Pi_e(\mathcal S_1+\mathcal S_2)\bigr|
 \le2\Pi_e\mathcal T_e.
\end{equation}
\end{lemma}

\begin{proof}
Split each of the two series of Lemma~\ref{an:balanced-afe} at $N_e$. By
Lemma~\ref{an:binning}, each finite part differs from its approximation
$\mathcal S_1$ or $\mathcal S_2$ by at most $\sum_i\rho_{20}(\xi_i)A_i$; the
same bound serves both series, since their two functions have the same
envelope. By Lemma~\ref{an:tails}, each tail is at most
$\zeta(\omega)^2(N_e+1)^{\omega+\nu}\mathcal V(t(N_e+1))$. Multiplying by
$\Pi_e$ gives \eqref{an:row-enclosure}.
\end{proof}

\begin{proposition}\label{an:genus-value}
For $\sigma=301/300$,
\[
 0.08264460177456<\frac{\log\zeta_E(\sigma)}{512}<0.08264460807138.
\]
\end{proposition}

\begin{proof}[Proof (computer-assisted)]
By Proposition~\ref{gf:factorization},
\[
 \log\zeta_E(\sigma)=\log\zeta(\sigma)+\log L(\sigma,\chi_B)
 +\sum_{e\ne0}\log L(\sigma,\chi_e).
\]
The coefficients $a_n(e)$ for $n\le N_e$ are computed exactly from the Euler
factors of Proposition~\ref{gf:local-data} by the supplementary program
\texttt{lfun241.py}. The supplementary program \texttt{afe241.py} evaluates
\eqref{an:row-enclosure} in Arb midpoint--radius interval arithmetic
\cite{Johansson2017,FLINT} with $256$ bits of precision, rounding every
quantity outward. Its routines for
Lemmas~\ref{an:balanced-afe}--\ref{an:binning} are taken unchanged from an
earlier supplementary archive of the author (Subsection~\ref{intro:code}), and
they implement these lemmas as stated here. Every enclosure of
$L(\sigma,\chi_e)$ is positive. The factor
$\zeta_B(\sigma)=\zeta(\sigma)L(\sigma,\chi_B)$ is evaluated with
\[
 L(\sigma,\chi_B)=241^{-\sigma}\sum_{a=1}^{240}\Bigl(\frac
 a{241}\Bigr)\zeta\Bigl(\sigma,\frac a{241}\Bigr),
\]
where $\zeta(s,x)=\sum_{n\ge0}(n+x)^{-s}$ is the Hurwitz zeta function,
using Arb's rigorous implementation of it \cite{Johansson2015Hurwitz}. This
gives $\zeta_B(\sigma)=725.51409164864\ldots$ and
$L(\sigma,\chi_B)=2.41373420241\ldots$. Summing the logarithms gives the
enclosure
\[
 [0.08264460177456168\ldots,\ 0.08264460807137072\ldots]
\]
of $\log\zeta_E(\sigma)/512$, which the proposition states rounded outward.
The logarithms are summed in Arb before any output conversion. The JSON
endpoints for the individual factors are rounded outward to double
precision, and their serialized decimal values are checked against the
Arb enclosures. The bounds for the normalized logarithm are retained as
Arb midpoint--radius strings, including their radii.
So the inequalities follow from
Lemmas~\ref{an:balanced-afe}--\ref{an:row}, the exact coefficients, and
outward-rounded interval arithmetic.
\end{proof}

As a check
independent of the approximate functional equation, the supplementary
programs \texttt{yecheck.gp} and \texttt{lvalues255.gp} evaluate the same
values numerically with PARI/GP \cite{PARI}: each of the $255$ values
$L(\sigma,\chi_e)$ lies in its enclosure, and PARI's value
$0.08264460492293\ldots$ of $\log\zeta_E(\sigma)/512$ lies in the enclosure
above.

\subsection{Frobenius Vectors and Lower Bounds for the Types}
\label{an:census}

\begin{lemma}\label{an:sigma2}
The set $\Sigma_2$ of Definition~\ref{an:P4} has exactly $21$ elements: zero
and the following twenty vectors, written as strings $c_0c_1\cdots c_7$ of
their coordinates, as in Section~\ref{tw:tower}:
\[
\begin{array}{ccccc}
 00100000&11110010&11010010&00010000&10000001\\
 10010001&00001000&10100111&10101111&00000100\\
 11101011&11101111&00000010&01100101&01100111\\
 00000001&01011010&01011011&10111010&11000101.
\end{array}
\]
In the notation of Table~\ref{tw:vector-table}, these twenty vectors are the
elementary images of the local involutions $x_j$, $y_j$ and $x_jy_j$ for
$j=1,2$, of the nonidentity elements of the four local groups
$C_2\times C_2$ above $3$ and $5$, and of $\iota_1$ and $\iota_2$.
\end{lemma}

\begin{proof}
Each of these images lies in $\Sigma_2$ by Lemma~\ref{tw:quadratic-layer},
since an involution has trivial square. The supplementary program
\texttt{census241.py} computes $\Sigma_2$ by linear algebra over $\Ftwo$
from the relations of Lemma~\ref{tw:quadratic-layer}, and finds no other
vectors.
\end{proof}

\begin{proposition}[Frobenius vectors of the primes of norm at most $10^6$]\label{an:census-prop}
Among the $78616$ primes
$\mathfrak p\notin S\cup\{\mathfrak t_0,\mathfrak t_1,\mathfrak t_2\}$ of
$B$ with $\mathrm N\mathfrak p\le10^6$, exactly $246$ have
$\overline{\Frob}_{\mathfrak p}=0$, exactly $5994$ have
$\overline{\Frob}_{\mathfrak p}\in\Sigma_2\setminus\{0\}$, and the remaining $72376$ form
$\mathcal P_4$.
\end{proposition}

\begin{proof}
The supplementary program \texttt{census241.py} computes the Frobenius
vector of each of these primes from Proposition~\ref{gf:local-data}(1) with
exact modular arithmetic, and compares it with the elements of $\Sigma_2$
(Lemma~\ref{an:sigma2}).
\end{proof}

Below we use this classification of the primes, not the three counts.

\begin{lemma}\label{an:corrections}
For $\sigma=301/300$,
\begin{equation}\label{an:census-values}
 \Delta_{2,7,29}(\sigma)>0.0344534299,\qquad
 \Delta_{\mathcal P_4}(\sigma)>0.0016298411.
\end{equation}
\end{lemma}

\begin{proof}
The supplementary program \texttt{ceiling241.py} uses the classification of
Proposition~\ref{an:census-prop} and evaluates \eqref{an:sel} and
\eqref{an:cen} in Arb interval arithmetic. It gives
$\Delta_{2,7,29}(\sigma)=0.03445342991613\ldots$ and
$\Delta_{\mathcal P_4}(\sigma)=0.00162984111205\ldots$, and the lemma states
these values rounded down.
\end{proof}

For every prime $\mathfrak p$ of $B$ put
\begin{equation}\label{an:floors}
 (e^0_{\mathfrak p},f^0_{\mathfrak p})=
 \begin{cases}
 (8,4)&\mathfrak p=\mathfrak p_1,\mathfrak p_2,\\
 (2,2)&\mathfrak p=\mathfrak q_1,\mathfrak q_2,\mathfrak r_1,\mathfrak r_2,\\
 (1,4)&\mathfrak p\in\{\mathfrak t_0,\mathfrak t_1,\mathfrak t_2\}\cup\mathcal P_4,\\
 (1,2)&\mathrm N\mathfrak p\le10^6,\ \overline{\Frob}_{\mathfrak p}\in\Sigma_2\setminus\{0\},\\
 (1,1)&\text{otherwise},
 \end{cases}
\end{equation}
where $\mathfrak q_1,\mathfrak q_2$ and $\mathfrak r_1,\mathfrak r_2$ are the
primes above $3$ and $5$ (Lemma~\ref{tw:base}).

\begin{lemma}\label{an:floors-valid}
For every $K\in\mathcal K$ and every prime $\mathfrak p$ of $B$,
$e_K(\mathfrak p)\ge e^0_{\mathfrak p}$ and
$f_K(\mathfrak p)\ge f^0_{\mathfrak p}$. Thus \eqref{an:floors} gives lower
bounds for the types in $K/B$, as Proposition~\ref{an:prime-budget}
requires.
\end{lemma}

\begin{proof}
At $\mathfrak p_j$ and $\mathfrak t_k$ the types in $K/B$ are $(8,4)$ and
$(1,4)$ by Theorem~\ref{tw:field-family}, as in the proof of
Proposition~\ref{an:Ystar}. At $\mathfrak q_j$ and $\mathfrak r_j$ the type
in $E/B$ is $(2,2)$ by Corollary~\ref{gf:types}, and if
$\overline{\Frob}_{\mathfrak p}\ne0$, then $f_E(\mathfrak p)=2$ by the same corollary.
Since $E\subseteq K$, the ramification indices and residue degrees in $E/B$
divide those in $K/B$ (Lemma~\ref{an:local-contribution}). On
$\mathcal P_4$ use Lemma~\ref{an:census-degree}. The bounds $(1,1)$ hold
trivially.
\end{proof}

\begin{lemma}\label{an:D}
For the lower bounds \eqref{an:floors},
\begin{equation}\label{an:R}
 \sum_{\mathfrak p}\frac{\log\mathrm N\mathfrak p}
 {2e^0_{\mathfrak p}(\mathrm N\mathfrak p^{2f^0_{\mathfrak p}}-1)}
 <0.0085105513.
\end{equation}
Hence $\mathcal D=0.0085105513$ satisfies the hypothesis of
Proposition~\ref{an:prime-budget}.
\end{lemma}

\begin{proof}
The primes of norm greater than $X=10^6$ receive $(1,1)$. At most two primes
of $B$ have a given norm, and $x\mapsto\log x/x^2$ decreases on
$[2,\infty)$. Their contribution to the sum is therefore at most
\[
 \sum_{m>X}\frac{\log m}{m^2-1}
 \le\frac{1}{1-X^{-2}}\sum_{m>X}\frac{\log m}{m^2}
 \le\frac{1}{1-X^{-2}}\int_X^\infty\frac{\log x}{x^2}\,dx
 =\frac{\log X+1}{X(1-X^{-2})}.
\]
The supplementary program \texttt{ceiling241.py} sums the primes of norm at
most $X$ in Arb interval arithmetic and evaluates this tail bound entirely
in Arb, including the logarithm. The resulting upper endpoint is less
than $0.0085105513$.
\end{proof}

\subsection{\texorpdfstring{Proof of Theorem~\ref{an:ceiling}}{Proof of the Upper Bound}}

\begin{corollary}\label{an:Ystar-value}
Let $\sigma=301/300$. Every $K\in\mathcal K$ satisfies
$[K:\Q]^{-1}\log\zeta_K(\sigma)\le Y_*$, where
\[
 Y_*=0.0826446081-0.0344534299-0.0016298411=0.0465613371.
\]
\end{corollary}

\begin{proof}
Combine Proposition~\ref{an:Ystar}, the upper bound of
Proposition~\ref{an:genus-value} rounded up, and the lower bounds of
Lemma~\ref{an:corrections}, which are rounded down.
\end{proof}

\begin{proof}[Proof of Theorem~\ref{an:ceiling}]
Let $\sigma=301/300$ and $\varepsilon=1/300$. We apply
Proposition~\ref{an:prime-budget} with $Y_*=0.0465613371$
(Corollary~\ref{an:Ystar-value}), with the lower bounds \eqref{an:floors}
for the types (Lemma~\ref{an:floors-valid}), and with
$\mathcal D=0.0085105513$ (Lemma~\ref{an:D}). Numerically
$\kappa_\infty=(\ell-\gamma-\log(4\pi))/4=0.63694120292\ldots<0.6369412030$,
so
\[
 Y_*+\varepsilon(\kappa_\infty+\mathcal D)
 <0.0465613371+\frac{0.6369412030+0.0085105513}{300}
 <0.0487128430<C.
\]
The interval evaluation of the same expression from the unrounded
enclosures, by \texttt{ceiling241.py}, gives the upper bound
$0.0487128429$. The replay compares the complete serialized Arb enclosure,
including its printed radius, with $C$. Proposition~\ref{an:prime-budget} with $C'=C$ proves
the theorem.
\end{proof}

\begin{remark}\label{an:limit}
The primes above $2$, $3$, $5$, $7$ and $29$ alone limit what
Proposition~\ref{an:prime-budget} can give. Let $(K^{(j)})$, $\beta_q$ and
$Z(s)=\sum_q\beta_qg_s(q)$ be as in the proof of that proposition. The
absolute types of the primes of $K^{(j)}$ above $2,3,5,7,29$ are known exactly
(Theorem~\ref{tw:field-family}), so these primes contribute exactly
\[
 \frac{2g_1(2^4)}{64}+\frac{2g_1(3^2)}8+\frac{2g_1(5^2)}8
 +\frac{2g_1(29^4)}8+\frac{g_1(7^8)}8=0.0416684615\ldots
\]
to $Z(1)$, a value printed by \texttt{ceiling241.py}; all other
contributions are nonnegative. The proof of
Proposition~\ref{an:prime-budget} shows that
$Z(1)\le Y_*+\varepsilon(\kappa_\infty+\mathcal D)$ for every choice of
$\sigma$, $Y_*$, $\mathcal D$ and lower bounds for the types that satisfies
its hypotheses. So no such choice gives an upper bound below this number.
The bound $C=\ceiling$ exceeds it by $0.0070444$.
\end{remark}

\begin{remark}[The hypothesis of the Lean formalization]\label{an:lean-hypothesis}
The Lean formalization \cite{NaslundLean2026} proves, under one explicit
hypothesis that it does not prove, that there are finite planar sets $U_j$
with $|U_j|\to\infty$ and $u(U_j)/|U_j|^{1.0427}\to\infty$. The hypothesis is
\begin{equation}\label{an:h241}
 \frac{\log\zeta_E(1+\frac1{300})}{512}
 +\frac1{300}\Bigl(\frac{\ell-\gamma-\log(4\pi)}4
 -\frac1{512}\frac{\zeta_E'}{\zeta_E}(2)\Bigr)<0.0852.
\end{equation}
In the formalization, $E$ is the subfield of an algebraic closure of $\Q$
generated by $\sqrt{241}$ and square roots of $\alpha_0,\ldots,\alpha_7$; it
is isomorphic to the Kummer field $E$, and the formalization calls it the
genus field.

The left side of \eqref{an:h241} is the bound of this section with all
local data taken from the fixed subfield $E$. By \eqref{an:local-debit} for
$M=E$, $-(\zeta_E'/\zeta_E)(2)/512$ is the sum defining $\mathcal D$ in
Proposition~\ref{an:prime-budget} when the lower bounds are the types in
$E/B$. These lower bounds are valid since $E\subseteq K$. The choice
$Y_*=\log\zeta_E(\sigma)/512$ is valid by Lemma~\ref{an:local-contribution}.
So the left side of \eqref{an:h241} is the bound given by
Proposition~\ref{an:prime-budget} without the refinements of
Proposition~\ref{an:Ystar} and with the types in $E/B$ in place of those in
$K/B$. The supplementary program \texttt{h241\_receipt.py} bounds it by
$0.0848335193<0.0852$. It uses the upper endpoint in
Proposition~\ref{an:genus-value}, the value of $\kappa_\infty$, and
$-(\zeta_E'/\zeta_E)(2)/512<0.0197321539$, which it obtains by summing
\eqref{an:local-debit} in interval arithmetic over the primes of norm at
most $10^6$ and bounding the rest as in the proof of Lemma~\ref{an:D}.

From \eqref{an:h241}, the formalization subtracts corrections at the primes
above $2$, $7$ and $29$ and at the primes above $41$, $47$, $53$, $59$,
$61$, $67$, $79$, $83$ and $97$, which lie in $\mathcal P_4$. It obtains the
weaker bound $0.0495$ for $d^{-1}\log L_F(1)$ for the fields of the tower
in the formalization. This suffices for the exponent $1.0427$ but not for
$1.04273$. The bound $C=\ceiling$ of Theorem~\ref{an:ceiling} uses
Proposition~\ref{an:census-prop} and the lower bounds for the types in
$K/B$; this refinement is not formalized.
\end{remark}
